\documentclass[10pt]{article}
\usepackage[a4paper,margin=1.5cm]{geometry}
\usepackage{amsmath}
\usepackage{xcolor}
\usepackage{enumitem}
\usepackage{titlesec}
\usepackage{parskip}
\usepackage[hidelinks]{hyperref}

\definecolor{accent}{RGB}{20,80,140}
\titleformat{\section}{\normalfont\bfseries\color{accent}\large}{}{0em}{}
\titlespacing{\section}{0pt}{5pt}{2pt}
\setlist[itemize]{leftmargin=1.2em,itemsep=1pt,topsep=1pt,parsep=0pt}
\pagestyle{empty}
\newcommand{\code}[1]{\texttt{\small #1}}

\begin{document}

\begin{center}
  {\LARGE\bfseries Implementing a 1D PRD source function}\\[2pt]
  {\large\color{accent} Full wavelength- and height-dependent $S(\lambda,z)$ applied to 3D opacity}\\[2pt]
  {\small Implementation plan \textperiodcentered\ 2026-09-08}
\end{center}

\section{Idea}
For a scattering resonance line, \emph{opacity is local} (keep computing it from
the real 3D $T$, $v$, populations) but the \emph{source function is non-local};
borrow the full PRD $S(\lambda,z)$ from the 1D NLTE model and integrate the
\emph{real} formal solution:
\[
  I_\lambda=\sum_i\bigl(1-e^{-\Delta\tau_{\lambda,i}}\bigr)e^{-\tau_{\lambda,<i}}\,
  S(\lambda',h_i'),\qquad
  \eta_i(\lambda)=\chi_i(\lambda)\,S(\lambda',h_i').
\]
A \textbf{single} source function is used --- \emph{no line/continuum split}. Off
limb the continuum contributes a negligible pedestal, and near line core
$\chi_{\rm line}\!\gg\!\chi_{\rm cont}$, so the total $S=\eta/\chi$ from the file is
already the line value there; applying it to the total 3D opacity is accurate
across the window and removes the continuum-isolation step entirely.

The curved ray already yields the geometric height $h_i$ of every sample point,
and the tangent height sets the minimum, so $S$ is sampled over $h\ge
h_{\rm tangent}$ --- precisely the off-limb regime. The lookup coordinates are
velocity- and calibration-shifted: $\lambda'=\lambda-\Delta\lambda_D(v_i)$,
$h_i'=h_i-\Delta z_{\rm cal}$ (\S4--5).

\section{Build the PRD table --- once, on rank 0}
\begin{itemize}
  \item Load \code{disk\_center\_test\_opem.fits} $\Rightarrow$ $\chi(z,\lambda)$,
    $\eta(z,\lambda)$, $z_{\rm 1D}$ (82 pts), $\lambda_{\rm 1D}$ (1001).
    Interpolate in $\lambda$ onto the synthesis grid \code{wave} so the two grids
    are guaranteed aligned.
  \item Form the total source function directly:
    $S(\lambda,z)=\eta(\lambda,z)/\chi(\lambda,z)$ --- no subtraction, no guard
    ($\chi$ never vanishes thanks to the continuum floor).
  \item Result: a compact table \code{\{lam\_1d, z\_1d, S[N$\lambda$,Nz]\}}
    ($\sim$0.66\,MB).
\end{itemize}

\section{Distribute to workers via MPI broadcast \small(adjustment 1)}
Only rank 0 reads the file. Before the manager/worker loop, all ranks call once
\[
  \code{s1d = comm.bcast(s1d, root=0)}
\]
(rank 0 passes the table, others pass \code{None}); each worker holds it for the
whole run. This costs one transfer per worker ($\sim$0.66\,MB, negligible) and
\textbf{no worker opens the file on disk}. It is independent of the existing
tag-based \code{send}/\code{recv} task machinery.

\section{Apply $S(\lambda,z)$ in the LOS integration \small(adjustment 2)}
\begin{itemize}
  \item \code{create\_ray}: also carry the height --- add \code{'Height': z} (km)
    to \code{param\_ray} (already computed, currently discarded).
  \item \code{calc\_op\_em}: compute the total opacity $\chi(\lambda,z)$ as now
    (line from pops $+$ LTE continuum). Build the lookup coordinates per sample
    point, $\lambda'_j=\lambda_j-\Delta\lambda_D(v_i)$ (same
    $\Delta\lambda_D=(v_i/c)\lambda_0$ already used for the profile $\varphi$) and
    $h'_i=h_i-\Delta z_{\rm cal}$; evaluate
    \code{RegularGridInterpolator((lam\_1d, z\_1d), S)} at all $(\lambda'_j,h'_i)$
    $\Rightarrow$ \code{S\_ray[N$\lambda$,Ns]}. Then set
    $\eta=\chi\cdot\code{S\_ray}$ (the population-based emissivity is no longer
    used for $S$). Return $\chi,\eta$.
  \item \code{synth}: call \code{simple\_formal\_solution(op, em, ds)} and
    \textbf{return the real $I$} instead of $1-e^{-\tau}$; keep the old behaviour
    behind a flag for side-by-side comparison.
  \item \code{worker\_work(rank, s1d)} / \code{synth(\dots, s1d)}: thread the
    broadcast table through. MPI overseer/worker structure is otherwise unchanged.
\end{itemize}

\section{Height calibration --- free parameter \small(adjustment 3)}
$\Delta z_{\rm cal}$ is a single scalar aligning MURaM's $z=0$ ($\tau_{500}{=}1$)
with FALC's $z=0$, with convention $z_{\rm 1D}^{\rm lookup}=h_{\rm ray}-\Delta
z_{\rm cal}$. Exposed as a top-level CLI argument (\code{sys.argv}), default $0$,
and applied \emph{at lookup time in the worker} --- so the broadcast table stays
offset-agnostic and one run can sweep $\Delta z_{\rm cal}$ without rebuilding it.
Calibrate by matching the disk-centre profile or the limb-intensity inflection.

\section{Caveats}
\begin{itemize}
  \item FALC tops out at \textbf{2.34\,Mm}; set \code{bounds\_error=False} and fill
    above the top with $S$ at the top (or 0). Tall spicules reach higher --- an
    extended/semi-empirical 1D model may be needed for the high LOS.
  \item The velocity shift assumes $S$ co-moves with the plasma (PRD in the atom
    frame); it can be disabled for a rest-frame first run.
  \item Line/continuum split is deliberately neglected here (justified off limb);
    revisit only if a wing or on-disk comparison later needs it.
  \item The broken \code{take\_given\_S} branch (undefined \code{otherids}) is
    replaced by this path.
\end{itemize}

\begin{center}\small\itshape
Scope: \code{create\_ray}, \code{calc\_op\_em}, \code{synth}, \code{worker\_work},
and \code{\_\_main\_\_} (load + \code{bcast} + new $\Delta z_{\rm cal}$ CLI arg).
\end{center}

\end{document}
